\documentclass[11pt]{letter}
\usepackage[margin=1in]{geometry}
\signature{Albert Linda}
\address{Department of Materials Science and Engineering\\ Indian Institute of Technology Kanpur\\ Kanpur 208016, Uttar Pradesh, India}
\begin{document}
\begin{letter}{Editors\\ Physical Review B}
\opening{Dear Editors,}

I submit the manuscript ``Local-resolvent formulation and assessment of linear-scaling Kohn--Sham calculations'' for
consideration in Physical Review B.

The manuscript writes the Kohn--Sham density in terms of the diagonal of the one-electron resolvent sampled at a few
complex energies and studies one simple realization of this idea. A half-step trapezoid rule on a circle
gives an occupation function that always lies between 0 and 1. We prove an a priori error bound for this rule and show that, with a circle that moves with the chemical
potential, forces at fixed electron number are exact at every number of poles.

The manuscript also reports where this realization falls short. Gauss--Legendre and Zolotarev pole rules need fewer poles, and the cost of the half-step rule is set by the ratio of occupied band width
to gap. In three dimensions the method is a divide-and-conquer scheme. We compare it with divide-and-conquer using
Fermi--Dirac occupations, which turns out to be equally accurate, and with dense and warm-started iterative Kohn--Sham
solvers in the same basis. For gapped H$_2$ molecular solids it reproduces self-consistent Kohn--Sham energies to
$2.5\times10^{-7}$~Ha per electron up to 1024 electrons, and, because of a large prefactor, complete calculations are slower than single-threaded dense diagonalization below about 500 electrons (by up to a factor of six at 54 electrons), 2.0 times faster at 686 electrons and 3.5 times faster at 1024 electrons (three interleaved repeats on an idle machine; no comparison with optimized or parallel dense solvers). The cases in which the method is inaccurate (disorder, compression, small gaps and diffuse
basis sets) are measured, and the relation to Green's-function, pole-expansion and
divide-and-conquer methods is discussed in the text.

All analytic identities were checked by computer algebra, and the implementation is covered by automated tests. The
Supplemental Material contains the solver comparisons, timing, memory and thread data, the basis-set, margin and pole-rule
studies, window diagnostics and SCF convergence histories. The code and data are available at \texttt{REPO\_URL}. The
manuscript is not under consideration elsewhere.

\closing{Sincerely,}
\end{letter}
\end{document}
